% ============================================================
\chapter{Variational Inference}
\label{ch:variational}
% ============================================================

MCMC methods are powerful but slow: each sample depends on the previous one, convergence is difficult to diagnose, and scaling remains problematic. In the 2000s, an alternative emerged: \emph{variational inference}, which transforms the sampling problem into an \emph{optimization} problem. The idea, formalized by Michael Jordan, Tommi Jaakkola, and collaborators, is to search within a family of simple distributions $\mathcal{Q}$ for the one that best approximates the posterior, by minimizing the Kullback--Leibler divergence. David Blei and his students democratized the approach in 2017 with ADVI (Automatic Differentiation Variational Inference), making variational inference as accessible as a function call. Today, it powers variational autoencoders (VAEs) and deep generative models.

\begin{intuition}[title=\textbf{Central idea}]
Variational inference transforms the Bayesian integration problem
(computing the posterior) into an \emph{optimization} problem. Instead of
sampling as in MCMC, we search for the distribution $q^*(\theta)$ in a
parametric family that best approximates the true posterior $p(\theta \mid x)$.
\end{intuition}

% ------------------------------------------------------------
\section{Motivation and problem setup}
% ------------------------------------------------------------

In many Bayesian models, the posterior
\[
  p(\theta \mid x) = \frac{p(x \mid \theta)\,\pi(\theta)}{p(x)}
\]
is intractable because the normalizing constant
$p(x) = \int p(x \mid \theta)\,\pi(\theta)\,d\theta$ has no closed form.
MCMC methods provide asymptotically exact solutions but can be prohibitively
expensive in high dimensions or with large datasets.

\begin{remark}
Variational inference is often preferred over MCMC when:
\begin{itemize}
  \item the dataset is very large ($n > 10^5$),
  \item a fast approximate answer is acceptable,
  \item the model has many latent variables (e.g., LDA, VAE).
\end{itemize}
\end{remark}

% ------------------------------------------------------------
\section{KL divergence and the ELBO}
% ------------------------------------------------------------

\begin{definition}[KL divergence]
\index{KL divergence}
Let $q$ and $p$ be two densities on $\Theta$. The \textbf{Kullback--Leibler
divergence} from $q$ to $p$ is:
\[
  \mathrm{KL}(q \| p) = \int q(\theta) \ln \frac{q(\theta)}{p(\theta)}\,d\theta.
\]
It satisfies $\mathrm{KL}(q \| p) \geq 0$ with equality if and only if
$q = p$ almost everywhere.
\end{definition}

The variational objective is to minimize $\mathrm{KL}(q \| p(\cdot \mid x))$
over a family $\mathcal{Q}$. Decomposing:
\[
  \ln p(x) = \underbrace{\E_q\!\bigl[\ln p(x, \theta) - \ln q(\theta)\bigr]}_{\mathrm{ELBO}(q)}
              + \mathrm{KL}\bigl(q \| p(\cdot \mid x)\bigr).
\]

\begin{definition}[ELBO --- Evidence Lower Bound]
\index{ELBO}
The \textbf{ELBO} (Evidence Lower Bound) is defined as:
\[
  \mathcal{L}(q) = \E_q\!\bigl[\ln p(x, \theta)\bigr] - \E_q\!\bigl[\ln q(\theta)\bigr].
\]
Since $\mathrm{KL} \geq 0$, we have $\mathcal{L}(q) \leq \ln p(x)$ for all $q$.
\end{definition}

\begin{theorem}[Equivalence of objectives]
Maximizing the ELBO $\mathcal{L}(q)$ with respect to $q \in \mathcal{Q}$
is equivalent to minimizing $\mathrm{KL}(q \| p(\cdot \mid x))$.
\end{theorem}

\begin{attention}
KL divergence is not symmetric: $\mathrm{KL}(q \| p) \neq \mathrm{KL}(p \| q)$.
Variational inference uses $\mathrm{KL}(q \| p)$ (the ``reverse'' or ``exclusive'' direction),
which makes $q$ mode-seeking and tend to underestimate the spread of $p$. Expectation
propagation uses $\mathrm{KL}(p \| q)$ (the ``forward'' or ``inclusive'' direction),
which makes $q$ cover the full support of $p$.
\end{attention}

% ------------------------------------------------------------
\section{Mean-field approximation}
% ------------------------------------------------------------

\begin{definition}[Mean-field family]
\index{mean-field}
Partition $\theta = (\theta_1, \ldots, \theta_K)$ and define:
\[
  \mathcal{Q}_{\mathrm{MF}} = \Bigl\{ q : q(\theta) = \prod_{k=1}^K q_k(\theta_k) \Bigr\}.
\]
Each factor $q_k$ is free (nonparametric within the factorization).
\end{definition}

\begin{theorem}[Optimal mean-field update]
\label{thm:cavi_update}
Fixing all $q_j$ ($j \neq k$), the optimal factor $q_k^*$ satisfies:
\[
  \ln q_k^*(\theta_k) = \E_{q_{-k}}\!\bigl[\ln p(x, \theta)\bigr] + \mathrm{const},
\]
where $q_{-k} = \prod_{j \neq k} q_j(\theta_j)$.
\end{theorem}

\begin{remark}
For conjugate exponential family models, each $q_k^*$ belongs to the same
family as the full conditional prior, making the updates analytic.
\end{remark}

% ------------------------------------------------------------
\section{CAVI --- Coordinate Ascent Variational Inference}
% ------------------------------------------------------------

\begin{definition}[CAVI algorithm]
\index{CAVI}
The \textbf{CAVI} algorithm cyclically applies the update from
Theorem~\ref{thm:cavi_update} for $k = 1, \ldots, K$ and monitors
convergence via the ELBO.
\end{definition}

\begin{proposition}[Convergence of CAVI]
Each CAVI update increases (or maintains) the ELBO. The algorithm
converges to a local maximum of $\mathcal{L}(q)$.
\end{proposition}

\begin{example}
Consider a univariate Gaussian model: $x_i \mid \mu, \tau \sim \mathcal{N}(\mu, \tau^{-1})$,
with $\mu \sim \mathcal{N}(0, \sigma_0^2)$ and $\tau \sim \mathrm{Gamma}(a_0, b_0)$.
The mean-field approximation $q(\mu, \tau) = q_\mu(\mu)\,q_\tau(\tau)$ gives:
\begin{align}
  q_\mu^*(\mu) &= \mathcal{N}\!\left(\frac{n\bar{x}\,\E[\tau]}{n\E[\tau] + \sigma_0^{-2}},\;
                   \frac{1}{n\E[\tau] + \sigma_0^{-2}}\right), \\
  q_\tau^*(\tau) &= \mathrm{Gamma}\!\left(a_0 + \tfrac{n}{2},\;
                   b_0 + \tfrac{1}{2}\sum_{i=1}^n \E_\mu\!\bigl[(x_i - \mu)^2\bigr]\right).
\end{align}
The cross-expectations are iterated until convergence.
\end{example}

% ------------------------------------------------------------
\section{Stochastic Variational Inference (SVI)}
% ------------------------------------------------------------

CAVI requires a full pass over all data at each iteration. For large
datasets, we use \textbf{SVI} (Stochastic Variational Inference).

\begin{definition}[SVI]
\index{SVI}
Distinguish \emph{global} variables $\beta$ and \emph{local} variables $z_i$.
At each iteration:
\begin{enumerate}
  \item Sample a mini-batch $S \subset \{1, \ldots, n\}$.
  \item Update local variational parameters $q(z_i)$ for $i \in S$.
  \item Compute the natural gradient of the ELBO with respect to the
        global parameters and take a gradient step.
\end{enumerate}
\end{definition}

\begin{proposition}[Natural gradient]
For an exponential family model, the natural gradient of the ELBO with
respect to the natural parameters $\lambda$ of $q(\beta)$ is:
\[
  \hat{\nabla}_\lambda \mathcal{L} = \lambda_{\mathrm{prior}} + n \cdot
  \E_{q(z_S)}\!\bigl[t(\beta, x_S, z_S)\bigr] - \lambda,
\]
where $t$ is the sufficient statistic and $S$ the mini-batch.
\end{proposition}

\begin{minted}{python}
# SVI pseudo-code
for epoch in range(max_epochs):
    S = random_minibatch(data, batch_size)
    # Local update
    for i in S:
        q_z[i] = update_local(q_beta, x[i])
    # Natural gradient for global parameters
    grad = lambda_prior + (n / len(S)) * sufficient_stats(S) - lam
    lam = lam + rho * grad  # rho = Robbins-Monro step size
\end{minted}

% ------------------------------------------------------------
\section{Reparameterization trick}
% ------------------------------------------------------------

\begin{definition}[Reparameterization trick]
\index{reparameterization}
If $\theta \sim q_\phi(\theta)$ can be written as $\theta = g(\phi, \epsilon)$
with $\epsilon \sim p(\epsilon)$ independent of $\phi$, then:
\[
  \nabla_\phi \E_{q_\phi}\!\bigl[f(\theta)\bigr]
  = \E_{p(\epsilon)}\!\bigl[\nabla_\phi f(g(\phi, \epsilon))\bigr],
\]
which enables backpropagation through the sampling step.
\end{definition}

\begin{example}
For $q_\phi(\theta) = \mathcal{N}(\mu, \sigma^2)$, set
$\theta = \mu + \sigma \epsilon$ with $\epsilon \sim \mathcal{N}(0,1)$.
This is the foundation of \textbf{Variational Autoencoders} (VAEs).
\end{example}

% ------------------------------------------------------------
\section{Comparison with MCMC}
% ------------------------------------------------------------

\begin{center}
\begin{tabular}{lcc}
\hline
\textbf{Criterion} & \textbf{MCMC} & \textbf{VI} \\
\hline
Nature & Sampling & Optimization \\
Convergence & Exact (asymptotic) & Approximate \\
Speed & Slow for large $n$ & Fast, scalable \\
Diagnostics & Traces, $\hat{R}$ & ELBO \\
Multimodality & Possible (HMC, \ldots) & Difficult \\
Uncertainty & Faithful & Underestimated (forward KL) \\
\hline
\end{tabular}
\end{center}

\begin{remark}
In practice, it is recommended to use VI for model exploration and MCMC
for final inference when accuracy is critical.
\end{remark}

% ------------------------------------------------------------
\section{Extensions: normalizing flows}
% ------------------------------------------------------------

One can enrich the variational family $\mathcal{Q}$ by composing
invertible transformations:
\[
  \theta_K = f_K \circ f_{K-1} \circ \cdots \circ f_1(\theta_0),
  \quad \theta_0 \sim q_0(\theta_0).
\]
The resulting density is obtained by the change-of-variables formula:
\[
  \ln q_K(\theta_K) = \ln q_0(\theta_0)
  - \sum_{k=1}^K \ln \abs{\det \frac{\partial f_k}{\partial \theta_{k-1}}}.
\]

% ------------------------------------------------------------
\section{Key formulas}
% ------------------------------------------------------------

\begin{keyformulas}
\begin{align}
  \text{ELBO:} \quad & \mathcal{L}(q) = \E_q[\ln p(x,\theta)] - \E_q[\ln q(\theta)] \\
  \text{Decomposition:} \quad & \ln p(x) = \mathcal{L}(q) + \mathrm{KL}(q \| p(\cdot \mid x)) \\
  \text{Mean field:} \quad & \ln q_k^*(\theta_k) = \E_{q_{-k}}[\ln p(x,\theta)] + C \\
  \text{Reparam.:} \quad & \nabla_\phi \E_{q_\phi}[f(\theta)]
    = \E_\epsilon[\nabla_\phi f(g(\phi,\epsilon))]
\end{align}
\end{keyformulas}

% ------------------------------------------------------------
\section{Exercises}
% ------------------------------------------------------------

\begin{exercise}
Show that $\mathrm{KL}(q \| p) \geq 0$ using Jensen's inequality.
Deduce that the ELBO is indeed a lower bound on $\ln p(x)$.
\end{exercise}

\begin{exercise}
Consider the model: $x_i \mid \mu \sim \mathcal{N}(\mu, 1)$, $i=1,\ldots,n$,
with $\mu \sim \mathcal{N}(0, \sigma_0^2)$. Compute the ELBO for
$q(\mu) = \mathcal{N}(m, s^2)$ and maximize it analytically in $m$ and $s^2$.
Verify that you recover the exact posterior.
\end{exercise}

\begin{exercise}
Implement the CAVI algorithm for the Gaussian model with unknown precision
from the example above. Plot the ELBO evolution across iterations.
\end{exercise}

\begin{exercise}
Experimentally compare (in Python) the mean-field variational approximation
and the Gibbs sampler for a two-component Gaussian mixture. Discuss the
observed differences in uncertainty estimation.
\end{exercise}

\begin{exercise}[ELBO computation for a Gaussian model]
Consider the model:
$x_1, \ldots, x_n \mid \mu, \sigma^2 \sim \mathcal{N}(\mu, \sigma^2)$
with $\sigma^2$ known, and the prior $\mu \sim \mathcal{N}(\mu_0, \tau_0^2)$.
Let $q(\mu) = \mathcal{N}(m, s^2)$ be the variational family.
\begin{enumerate}
  \item Write the ELBO $\mathcal{L}(m, s^2) =
        \E_q[\ln p(x_{1:n}, \mu)] - \E_q[\ln q(\mu)]$ explicitly as a
        function of $m$, $s^2$, $\mu_0$, $\tau_0^2$, $\sigma^2$,
        $\bar{x}$, and $n$.
  \item Maximize $\mathcal{L}$ with respect to $m$ and $s^2$ and show
        that the exact posterior $\mathcal{N}(m^*, (s^*)^2)$ is recovered.
  \item Evaluate the ELBO at the optimum and verify that
        $\mathcal{L}(m^*, (s^*)^2) = \ln p(x_{1:n})$ (the KL vanishes).
\end{enumerate}
\end{exercise}

\begin{exercise}[Mean-field update for a linear model]
Consider the Bayesian linear regression model:
$y \mid X, \beta, \tau \sim \mathcal{N}(X\beta, \tau^{-1} I_n)$,
$\beta \sim \mathcal{N}(0, \alpha^{-1} I_p)$,
$\tau \sim \mathrm{Gamma}(a_0, b_0)$, with $\alpha$ fixed.
Set $q(\beta, \tau) = q_\beta(\beta)\,q_\tau(\tau)$.
\begin{enumerate}
  \item Derive the optimal update $q_\beta^*(\beta)$ by computing
        $\E_{q_\tau}[\ln p(y, \beta, \tau)]$ and show that
        $q_\beta^*(\beta) = \mathcal{N}(\mu_\beta, \Sigma_\beta)$
        with $\Sigma_\beta = (\E[\tau]\,X^\top X + \alpha I_p)^{-1}$
        and $\mu_\beta = \E[\tau]\,\Sigma_\beta\,X^\top y$.
  \item Derive $q_\tau^*(\tau) = \mathrm{Gamma}(a_n, b_n)$ by
        computing $\E_{q_\beta}[\ln p(y, \beta, \tau)]$.
        Express $a_n$ and $b_n$.
  \item Implement the resulting CAVI algorithm and verify ELBO
        convergence on simulated data.
\end{enumerate}
\end{exercise}

\begin{exercise}[Comparison of VI quality metrics]
Let $p(\theta \mid x)$ be a bimodal posterior (mixture of two Gaussians
in one dimension) and $q(\theta) = \mathcal{N}(m, s^2)$.
\begin{enumerate}
  \item Numerically compute $\mathrm{KL}(q \| p)$ (reverse KL) and
        $\mathrm{KL}(p \| q)$ (forward KL) for various values of $m$
        and $s$. Graphically illustrate the respective minimizers.
  \item Show that the minimizer of $\mathrm{KL}(q \| p)$ concentrates
        on a single mode (``mode-seeking''), while the minimizer of
        $\mathrm{KL}(p \| q)$ covers both modes (``mass-covering'').
  \item Also compute the chi-squared divergence
        $\chi^2(q \| p) = \int \frac{(q - p)^2}{p}\,d\theta$ and the
        Wasserstein distance $W_2(q, p)$. Compare the variational
        approximations obtained by minimizing these four metrics.
\end{enumerate}
\end{exercise}
